\documentclass[11pt,a4paper]{article}
\usepackage[utf8]{inputenc}
\usepackage[T1]{fontenc}
\usepackage[spanish]{babel}
\usepackage{geometry}
\geometry{margin=2.5cm}
\usepackage{xcolor}
\usepackage{fancyhdr}

\definecolor{wasaffblue}{RGB}{0,82,147}

\pagestyle{fancy}
\fancyhf{}
\fancyhead[C]{\small\textcolor{wasaffblue}{WASAFF Consulting}}
\fancyfoot[C]{\thepage}
\renewcommand{\headrulewidth}{0.4pt}

\begin{document}

\begin{flushright}
Santiago, 18 de mayo de 2026
\end{flushright}

\vspace{0.5cm}

Estimados señores de \textbf{Kaeser Compresores Chile SpA} y \textbf{CMPC Papeles Cordillera S.A.}:

\vspace{0.5cm}

Por medio de la presente, me dirijo a ustedes en mi calidad de Director Técnico de WASAFF Consulting, para presentar formalmente nuestra propuesta técnica y comercial para el desarrollo de la \textbf{Ingeniería de Detalle del Sistema de Recuperación de Calor del Aceite Lubricante de Compresores Kaeser} en la planta de CMPC Papeles Cordillera, ubicada en Puente Alto.

\textbf{Propuesta de valor diferenciada:}

A diferencia de los enfoques convencionales basados en catálogos y factores de seguridad empíricos, nuestra propuesta integra \textbf{modelación matemática avanzada de nivel de maestría}, que incluye:

\begin{itemize}
  \item Dimensionamiento preciso del intercambiador mediante el método $\varepsilon$-NTU con propiedades termofísicas dependientes de temperatura.
  \item Análisis dinámico transitorio mediante integración Runge-Kutta de orden 4--5 (RK45), para cuantificar el comportamiento ante arranques, paradas y fallas.
  \item Análisis de sensibilidad paramétrica con tornado diagrams, identificando los parámetros críticos que impactan el dimensionamiento.
  \item Evaluación de integración con el sistema de cogeneración existente (Rolls-Royce Trent 60 + HRSG).
\end{itemize}

\textbf{Experiencia directa con CMPC Papeles Cordillera:}

Cuento con experiencia previa en el diseño de sistemas de recuperación de calor para esta misma planta, lo que me proporciona un conocimiento único de sus procesos, infraestructura y requisitos específicos. Adicionalmente, mi formación como Físico Computacional y mi actual cursada del Magíster en Simulación Computacional (PUCV) garantizan un nivel técnico acorde a los estándares de un grupo como CMPC.

\textbf{Documentos adjuntos:}

\begin{enumerate}
  \item Propuesta Técnico-Comercial WASAFF — CMPC Cordillera (11 páginas)
  \item Checklist de Información Requerida (formato CSV/Excel)
  \item Acuerdo de Confidencialidad (NDA) para firma
\end{enumerate}

\textbf{Próximos pasos sugeridos:}

\begin{enumerate}
  \item Revisión interna de la propuesta por parte de Kaeser y CMPC (5 días hábiles).
  \item Reunión técnica conjunta para validar alcance y resolver dudas.
  \item Entrega de información base por parte de CMPC (checklist adjunto).
  \item Ajuste de propuesta si es necesario y firma de contrato.
\end{enumerate}

Quedamos atentos a sus comentarios y disponibles para coordinar una reunión presencial o virtual la próxima semana.

\vspace{1cm}

\begin{flushright}
\textit{Atentamente,}

\vspace{1cm}

\textbf{Felipe Wasaff}\\
Físico Computacional, MSc Candidate\\
Director Técnico, WASAFF Consulting\\
Tel: +56 9 4612 5682\\
Email: felipe.wasaff@uchile.cl\\
Web: www.wasaff.cl
\end{flushright}

\vspace{1cm}

\begin{center}
\rule{0.8\textwidth}{0.5pt}
\end{center}

\begin{center}
\small\textit{WASAFF Consulting — Ingeniería Térmica y Modelación Computacional Avanzada}\\
\small Pasaje Desierto Florido Sur \#1328, Puente Alto, Santiago, Chile
\end{center}

\end{document}
